% preamble.tex --- shared by all logbook documentation files
% Engine: pdflatex

\documentclass[11pt, a4paper]{report}

% --- Encoding and fonts ---
\usepackage[T1]{fontenc}
\usepackage[utf8]{inputenc}
\usepackage{lmodern}

% --- Page geometry ---
\usepackage[
    top=2.5cm, bottom=2.5cm,
    left=2.8cm, right=2.8cm
]{geometry}

% --- Colours ---
\usepackage[dvipsnames]{xcolor}
\definecolor{pionorange}{RGB}{255,152,0}
\definecolor{codeblue}{RGB}{0,80,160}
\definecolor{notebg}{RGB}{255,248,230}
\definecolor{warningbg}{RGB}{255,235,235}
\definecolor{codebg}{RGB}{245,245,245}
\definecolor{tabhead}{RGB}{230,230,230}

% --- Hyperlinks ---
\usepackage[
    colorlinks=true,
    linkcolor=codeblue,
    urlcolor=codeblue,
    citecolor=codeblue,
    pdftitle={Logbook Entry Generator Documentation},
    pdfauthor={PionCT Experiment}
]{hyperref}

% --- Tables ---
\usepackage{booktabs}
\usepackage{array}
\usepackage{tabularx}
\usepackage{longtable}

% --- Code listings ---
\usepackage{listings}
\lstset{
    basicstyle=\small\ttfamily,
    backgroundcolor=\color{codebg},
    frame=single,
    framerule=0.4pt,
    rulecolor=\color{gray!40},
    breaklines=true,
    breakatwhitespace=true,
    showstringspaces=false,
    tabsize=2,
    xleftmargin=6pt,
    xrightmargin=6pt,
    aboveskip=8pt,
    belowskip=6pt,
}
\lstdefinelanguage{JSON}{
    basicstyle=\small\ttfamily,
    morestring=[b]",
    stringstyle=\color{codeblue},
    literate=
        *{:}{{{\color{black}{:}}}}{1}
         {\{}{{{\color{black}{\{}}}}{1}
         {\}}{{{\color{black}{\}}}}}{1}
         {[}{{{\color{black}{[}}}}{1}
         {]}{{{\color{black}{]}}}}{1},
}
\lstdefinelanguage{bash}{
    basicstyle=\small\ttfamily,
    morecomment=[l]{\#},
    commentstyle=\color{gray},
}
\lstdefinelanguage{javascript}{
    basicstyle=\small\ttfamily,
    keywords={const,let,var,function,return,if,else,for,await,async,true,false},
    keywordstyle=\color{codeblue},
    morestring=[b]',
    morestring=[b]",
    stringstyle=\color{BrickRed},
    morecomment=[l]{//},
    commentstyle=\color{gray},
}

% --- Inline code ---
\newcommand{\code}[1]{\texttt{\small#1}}

% --- Note / Warning boxes ---
\usepackage{tcolorbox}
\tcbuselibrary{skins, breakable}

\newtcolorbox{notebox}{
    enhanced, breakable,
    colback=notebg,
    colframe=pionorange,
    leftrule=4pt, rightrule=0pt, toprule=0pt, bottomrule=0pt,
    arc=0pt, outer arc=0pt,
    left=8pt, right=8pt, top=4pt, bottom=4pt,
    fontupper=\small,
}
\newtcolorbox{warningbox}{
    enhanced, breakable,
    colback=warningbg,
    colframe=red!70!black,
    leftrule=4pt, rightrule=0pt, toprule=0pt, bottomrule=0pt,
    arc=0pt, outer arc=0pt,
    left=8pt, right=8pt, top=4pt, bottom=4pt,
    fontupper=\small,
}

% --- Section styling ---
\usepackage{titlesec}
\titleformat{\chapter}[block]
    {\normalfont\Large\bfseries\color{pionorange}}
    {\thechapter.}{0.8em}{}
    [\vspace{-0.3em}\color{pionorange}\rule{\linewidth}{1.5pt}]
\titlespacing{\chapter}{0pt}{0pt}{12pt}

\titleformat{\section}
    {\normalfont\large\bfseries\color{black}}
    {\thesection}{0.8em}{}
\titleformat{\subsection}
    {\normalfont\normalsize\bfseries}
    {\thesubsection}{0.8em}{}

% --- Header / footer ---
\usepackage{fancyhdr}
\pagestyle{fancy}
\fancyhf{}
\fancyhead[L]{\small\color{gray}\leftmark}
\fancyhead[R]{\small\color{gray}Logbook Entry Generator}
\fancyfoot[C]{\small\thepage}
\renewcommand{\headrulewidth}{0.4pt}
\renewcommand{\chaptermark}[1]{\markboth{#1}{}}

% --- Misc ---
\usepackage{enumitem}
\setlist[itemize]{itemsep=2pt, topsep=4pt}
\setlist[enumerate]{itemsep=2pt, topsep=4pt}
\usepackage{parskip}
\setlength{\parskip}{6pt}
\setlength{\parindent}{0pt}


\begin{document}

\title{
    \vspace{2cm}
    {\color{pionorange}\rule{\linewidth}{2pt}}\\[0.6em]
    {\Huge\bfseries Logbook Entry Generator}\\[0.3em]
    {\Large Setup Guide}\\[0.4em]
    {\color{pionorange}\rule{\linewidth}{2pt}}\\[1em]
    {\large PionCT Experiment --- Jefferson Lab}
    \vspace{1cm}
}
\author{}
\date{\today}
\maketitle
\thispagestyle{empty}

\tableofcontents
\newpage

% ============================================================
\chapter{Introduction}
% ============================================================

This guide is for experiment coordinators who need to configure and distribute
the logbook tool for their experiment. It covers editing the configuration
file, running the build script, and setting up Google Sheets export.

% ============================================================
\chapter{Prerequisites}
% ============================================================

\begin{itemize}
    \item Python 3.6 or later (standard library only --- no \code{pip} installs
          needed)
    \item A terminal (any operating system)
    \item Optionally: a Google account with access to Google Drive/Sheets, if
          you want the Sheets export feature
\end{itemize}

% ============================================================
\chapter{File Layout}
% ============================================================

After downloading, your working directory should look like this:

\begin{lstlisting}
my_experiment/
+-- build.py          <- build script (do not edit)
+-- config.json       <- your experiment configuration (edit this)
+-- logo.png          <- your experiment logo (optional)
`-- output/           <- generated files appear here (created automatically)
    +-- logbook_generator.html
    `-- apps_script.gs   (only if Google Sheets is enabled)
\end{lstlisting}

Distribute \code{output/logbook\_generator.html} to your shift workers.
Everything they need is embedded in that single file --- no other files need to
travel with it. The logo is base64-embedded automatically if it is found next
to \code{config.json} at build time.

% ============================================================
\chapter{Editing \texttt{config.json}}
% ============================================================

Open \code{config.json} in any text editor. The following sections describe
each configurable field.

\section{Experiment Identity}

\begin{lstlisting}[language=JSON]
"experiment_name": "PionCT",
"page_title":      "PionCT Logbook Entry Generator",
"contact_email":   "gayoso@jlab.org",
"credit_line":     "Original idea from Andrew Schick for the PRadII experiment"
\end{lstlisting}

\begin{itemize}
    \item \code{experiment\_name} --- used to namespace the browser's
          localStorage key so sessions from different experiments do not
          collide. Keep it short with no spaces.
    \item \code{page\_title} --- displayed as the H1 heading and browser tab
          title.
    \item \code{contact\_email} --- shown at the bottom of the Session
          Management box.
    \item \code{credit\_line} --- optional acknowledgement line. Remove the
          key entirely to omit it.
\end{itemize}

\section{Logo}

\begin{lstlisting}[language=JSON]
"logo": {
  "path": "logo.png",
  "alt":  "My Experiment Logo"
}
\end{lstlisting}

\begin{itemize}
    \item \code{path} --- relative to \code{config.json}. If the file exists
          there, the build script embeds it as base64 in the HTML (fully
          self-contained, no path dependency). If it does not exist, the path
          is used as-is (suitable for hosted URLs such as \code{https://...}).
    \item Set \code{"logo": \{\}} or remove the key entirely to show no logo.
\end{itemize}

\section{Quick Links}

\begin{lstlisting}[language=JSON]
"quick_links": [
  { "label": "Short Term Run Plan",           "id": "runPlanLink"   },
  { "label": "Link to Shift Worker Checklist","id": "checklistLink" }
]
\end{lstlisting}

Each entry creates a URL input field in the Summary section. \code{label} is
the display text; \code{id} must be a unique alphanumeric identifier (no
spaces). Add, remove, or relabel as many as your experiment needs. The import
feature uses the label text to identify links when re-parsing a generated HTML
entry.

\section{Google Sheets}

\begin{lstlisting}[language=JSON]
"google_sheets": {
  "enabled":       true,
  "shared_secret": "CHANGE_ME_TO_SOMETHING_PRIVATE"
}
\end{lstlisting}

\begin{itemize}
    \item Set \code{"enabled": false} to remove all Sheets UI from the
          generated HTML and skip generating \code{apps\_script.gs}.
    \item \code{shared\_secret} --- any string you choose. It must match the
          value in the deployed Apps Script. Keep it private; it is the only
          access control on who can POST rows to your sheet.
\end{itemize}

\section{Run Fields}

This is the most important section. It defines every column in the run summary
table and (if Sheets is enabled) the Google Sheet.

\begin{lstlisting}[language=JSON]
"run_fields": [
  {
    "key":          "run-number",
    "label":        "Run Number",
    "type":         "text",
    "source":       "user",
    "sheet_column": "Run Number",
    "sheet_align":  "right"
  },
  ...
]
\end{lstlisting}

\begin{center}
\begin{tabularx}{\linewidth}{@{}llp{2.2cm}X@{}}
\toprule
\textbf{Property} & \textbf{Required} & \textbf{Values} & \textbf{Meaning} \\
\midrule
\code{key}          & yes & hyphenated & CSS class used internally. Must be unique. \\
\code{label}        & yes & any string & Label shown in the form and table header. \\
\code{type}         & yes & \code{text}, \code{textarea} & Input type. Use \code{textarea} for Comments. \\
\code{source}       & yes & see below  & Where the value comes from. \\
\code{sheet\_column}& yes & any string & Column header written to Google Sheets. \\
\code{sheet\_align} & yes & \code{left}, \code{right}, \code{center} & Alignment in the HTML table. \\
\code{placeholder}  & no  & any string & Placeholder text in the input. \\
\code{\_comment}    & no  & any string & Ignored by build script. For your notes. \\
\bottomrule
\end{tabularx}
\end{center}

\subsection*{\texttt{source} values}

\begin{itemize}
    \item \code{"user"} --- shown as an editable field in the form; sent to
          the sheet.
    \item \code{"auto"} --- shown read-only in the form (greyed out); computed
          automatically. Two auto fields are currently supported:
          \begin{itemize}
              \item Key \code{"date"} --- filled with the local date
                    (\code{YYYY-MM-DD}) at generate/send time. Not shown in
                    the form at all.
              \item Key \code{"run-length"} --- auto-calculated as Stop
                    $-$ Start in minutes. Shown in the form but read-only.
          \end{itemize}
    \item \code{"sheet-only"} --- not shown in the HTML form at all. Written
          as an empty cell to the sheet (to be filled manually in the
          spreadsheet later).
\end{itemize}

\begin{warningbox}
\textbf{Warning:} The two \code{auto} keys (\code{date} and \code{run-length})
have special hardcoded behaviour in the build script. Any other field with
\code{source: "auto"} will appear as a greyed-out text input but will not be
auto-populated.
\end{warningbox}

\textbf{Field order matters.} The order of entries in \code{run\_fields}
determines the left-to-right column order in both the HTML table and the
Google Sheet. See Chapter~\ref{ch:schema} if you change the order after
deploying the Apps Script.

% ============================================================
\chapter{Running the Build Script}
% ============================================================

\begin{lstlisting}[language=bash]
python3 build.py             # uses config.json in the current directory
python3 build.py my_exp.json # uses a specific config file
\end{lstlisting}

Output is written to \code{output/} next to the config file. The script will:

\begin{enumerate}
    \item Validate \code{config.json} and exit with a clear error if anything
          is wrong.
    \item Embed the logo as base64 if the file is found (prints a
          confirmation).
    \item Write \code{output/logbook\_generator.html}.
    \item Write \code{output/apps\_script.gs} (only if
          \code{google\_sheets.enabled} is true).
    \item Print the schema hash (see Chapter~\ref{ch:schema}).
\end{enumerate}

Rerun the script any time you change \code{config.json}. The output files are
always regenerated from scratch.

% ============================================================
\chapter{Schema Hash and Changing Fields Mid-Experiment}
\label{ch:schema}
% ============================================================

The build script computes a hash of your \code{run\_fields} column order and
saves it to a hidden file \code{.schema\_hash} next to \code{config.json}. On
the next build, it compares the new hash to the saved one.

If the hash changes and Sheets is enabled, the script prints:

\begin{lstlisting}
  +------------------------------------------------------------------+
  |  WARNING: SHEETS SCHEMA CHANGED                                  |
  |  The run field schema has changed since the last build.          |
  |  You MUST redeploy the Apps Script and update the Web App URL    |
  |  in the logbook tool's Google Sheets Export Settings, or the     |
  |  column order in your sheet will be wrong.                       |
  +------------------------------------------------------------------+
\end{lstlisting}

\textbf{What this means in practice:} if you add, remove, or reorder columns
in \code{run\_fields}, the existing Apps Script deployment maps columns by
position, not by name. A mismatch will silently write values into the wrong
columns.

\begin{notebox}
\textbf{Note:} If \code{.schema\_hash} is missing (e.g.\ it was deleted or
this is the first build), the comparison is skipped silently and a new hash
file is written. No data is lost.
\end{notebox}

The current schema hash is also embedded as a comment in the generated HTML
and at the top of \code{apps\_script.gs}, so you can always check which build
produced a given file.

% ============================================================
\chapter{Setting Up Google Sheets Export}
% ============================================================

\section{Step 1 --- Create the Google Sheet}

Create a new Google Sheet. The build script will write a header row
automatically the first time a row is appended, but you can also add the
header manually if you want to set column widths or formatting in advance. The
column order must match the \code{sheet\_column} values in \code{config.json}
in the same order as \code{run\_fields}.

\section{Step 2 --- Open Apps Script}

In the Google Sheet: \textbf{Extensions $\to$ Apps Script}.

Delete any placeholder code in the editor. Paste the entire contents of
\code{output/apps\_script.gs}.

\section{Step 3 --- Set the Shared Secret}

In the script, find this line near the top:

\begin{lstlisting}[language=javascript]
const SHARED_SECRET = 'CHANGE_ME_TO_SOMETHING_PRIVATE';
\end{lstlisting}

Replace the placeholder with whatever string you set in \code{config.json}
under \code{google\_sheets.shared\_secret}. They must match exactly.

\section{Step 4 --- Deploy as Web App}

Click \textbf{Deploy $\to$ New deployment}.

\begin{itemize}
    \item Type: \textbf{Web app}
    \item Execute as: \textbf{Me}
    \item Who has access: \textbf{Anyone}
\end{itemize}

Click \textbf{Deploy}, authorise the permissions when prompted, then copy the
\textbf{Web App URL} that appears. It looks like:\\
\code{https://script.google.com/macros/s/XXXXXXXXXX/exec}

\begin{warningbox}
\textbf{Security note:} ``Anyone'' means anyone with the URL can POST to the
script. The shared secret is the only access control. Keep the URL and secret
out of public repositories.
\end{warningbox}

\section{Step 5 --- Distribute to Shift Workers}

Open \code{output/logbook\_generator.html} in a browser. In the
\textbf{Google Sheets Export Settings} section, paste:

\begin{itemize}
    \item \textbf{Web App URL} --- the URL from Step 4
    \item \textbf{Shared Secret} --- the same string from your config
\end{itemize}

These are saved in the browser's localStorage automatically. Each shift worker
on a different machine needs to do this once.

\section{Redeploying After a Schema Change}

If you change \code{run\_fields} and the build script warns you about a schema
change:

\begin{enumerate}
    \item Rebuild to get the new \code{apps\_script.gs}.
    \item In Apps Script: \textbf{Deploy $\to$ Manage deployments $\to$ Edit
          (pencil icon)}.
    \item Change ``Version'' to \textbf{New version}.
    \item Click \textbf{Deploy}.
    \item The URL \textbf{does not change} on redeployment --- you do not need
          to update it in the logbook tool.
\end{enumerate}

% ============================================================
\chapter{Distributing the Tool}
% ============================================================

Send shift workers \code{output/logbook\_generator.html}. That is the only
file they need. The logo and all CSS are embedded. They open it locally in a
browser --- no web server required.

If you want to host it (e.g.\ on a lab web server or GitHub Pages), that works
too --- just upload the single HTML file.

\end{document}
